\documentclass[10pt,a4paper]{amsart}
\usepackage[margin=19mm]{geometry}
\usepackage{amsmath,amssymb,mathtools}
\usepackage[scr=rsfs,cal=euler]{mathalfa}
\usepackage{xfrac}
\usepackage[unicode,hidelinks,bookmarks=false]{hyperref}
\newcommand{\ie}{i.e.\ }
\newcommand{\cf}{cf.\ }
\newcommand{\resp}{respectively\ }
\newcommand{\Subsection}{\subsection}
\providecommand{\nobreakdash}{\nobreak\hbox{-}}
\setlength{\emergencystretch}{2em}
\multlinegap0pt
\usepackage{bm}
\usepackage{bbm}
\usepackage{dsfont}
\usepackage{xspace}
\usepackage{cancel}
\usepackage{mathtools}
\usepackage{amsthm}
\makeatletter
\@ifundefined{claims}{\newtheorem{claims}{Claims}}{}
\@ifundefined{lemma}{\newtheorem{lemma}{Lemma}}{}
\makeatother
\makeatletter
\newcommand{\C}{\mathbb{C}}
\newcommand{\HC}{\mathbb{H}_{\C}}
\expandafter\ifx\csname Mlist\endcsname\relax
\newenvironment{Mlist}{\begin{itemize}[topsep=0pt,itemsep=4pt,leftmargin=7mm]}{\end{itemize}}
\fi
\renewcommand{\P}{\mathbb{P}}
\newcommand{\cL}{\mathcal{L}}
\newcommand{\cN}{\mathcal{N}}
\newcommand{\cR}{\mathcal{R}}
\expandafter\ifx\csname claims\endcsname\relax
\newenvironment{claims}{\begin{enumerate}[topsep=0pt,itemsep=0pt,leftmargin=8mm,label=\textnormal{\textbf{(\alph*)}},ref=(\alph*)]}{\end{enumerate}}
\fi
\newcommand{\ii}{\mathfrak{i}}
\newcommand{\qi}{\mathbf{i}}
\newcommand{\qj}{\mathbf{j}}
\newcommand{\qk}{\mathbf{k}}
\newcommand{\set}[2]{\left\{#1 ~|~ #2\right\}}
\makeatother
\title[Bivariate quaternionic factorizations and surfaces that deco]{Bivariate quaternionic factorizations and surfaces that decompose into two circles}
\author{Johanna Frischauf, Niels Lubbes, Hans-Peter Schr\"ocker}
\date{Source: arXiv:2608.11887v1 (CC BY 4.0)}
\begin{document}
\maketitle

\noindent\emph{Source and licence.} Bivariate quaternionic factorizations and surfaces that decompose into two circles. Version arXiv:2608.11887v1 (CC BY 4.0), distributed under the Creative Commons Attribution 4.0 International licence, as recorded in the arXiv licence metadata for this exact version and shown on the abstract page https://arxiv.org/abs/2608.11887v1. Full rights notice: licenses/arxiv-2608.11887-CC-BY-4.0.txt.\par\smallskip

\noindent\textit{Editorial scope.} Counted unit: the Lemma whose source label is \texttt{lem:N} in the article (source0.tex lines 1139--1160, source label \texttt{lem:N}), delivered together with the article's own complete original proof of it (source0.tex lines 1162--1230, one \texttt{proof} environment of 69 lines). No mathematics is written, repaired or re-derived: statement and proof are the article's own text line for line. Required context retained uncounted: none. Every definition, display and label of those lines is retained unchanged. Omitted from this package and not redistributed: 1--1138, 1161--1161, 1231--3367. Nothing inside a retained excerpt is omitted or altered: every retained body is delivered exactly as the article's lines are written, so no omission receipt of any kind accompanies this package. Citations: the retained excerpts carry no citation command at all, so no bibliography entry of the article is transcribed and no reference is redistributed. Reference adaptation: none was needed and none was made; every reference inside the delivered unit resolves inside it, so no unresolved reference or citation is produced. Printed slips kept literal: no printed word, symbol, hypothesis or number of the counted statement or of its proof was altered. Layout and class: the article's own class is \texttt{article} with packages including amsmath, amssymb, amsthm, amsfonts, amscd, thmtools, natbib, hyperref, cleveref, xcolor, geometry, enumitem, fontenc, tikz; the wrapper is the lane's pinned \texttt{amsart} editorial wrapper at 10pt on A4, which restates the theorem-like environments the retained text needs and adds the article's own macro definitions that the retained text needs (\texttt{\textbackslash C}, \texttt{\textbackslash HC}, \texttt{\textbackslash P}, \texttt{\textbackslash cL}, \texttt{\textbackslash cN}, \texttt{\textbackslash cR}, \texttt{\textbackslash ii}, \texttt{\textbackslash qi}, \texttt{\textbackslash qj}, \texttt{\textbackslash qk}, \texttt{\textbackslash set}), with extra wrapper packages (bm, bbm, dsfont, xspace, cancel, mathtools). The delivered numbering is the unit's own. Assets: no retained span contains a figure environment, a graphics-inclusion command, a PDF-page inclusion, a table or a TikZ picture; no image file is copied into this package, the package declares no asset dependency and no image file is redistributed with it. Licence: version arXiv:2608.11887v1 of this article is distributed under the Creative Commons Attribution 4.0 International licence, as recorded in the arXiv licence metadata for this exact version and shown on the abstract page https://arxiv.org/abs/2608.11887v1. A full rights notice is delivered at licenses/arxiv-2608.11887-CC-BY-4.0.txt. Characters: every retained excerpt is pure ASCII, so the delivered engine, XeLaTeX with its default Latin Modern text font, reports no missing character.
\begin{lemma}
\label{lem:N}
Suppose that $\alpha,h,q\in\HC\setminus\{0\}$ such that $\alpha\cdot\alpha^*=0$.
\begin{claims}
\item\label{lem:N:a}
The generators $\cL_\alpha$ and $\cR_\alpha$ are the two complex lines in the hyperquadric $\cN$
that pass through the complex point $[\alpha^*]\in\cN$.

\item\label{lem:N:b}
If $[h],[q]\in \cL_\alpha$, then $h\cdot q^*=0$.\\
If $[h],[q]\in \cR_\alpha$, then $q^*\cdot h=0$.

\item\label{lem:N:c}
If $h\cdot q^*=h\cdot h^*=q\cdot q^*=0$, then $[h],[q]\in \cL_{q^*}$.\\
If $q^*\cdot h=h\cdot h^*=q\cdot q^*=0$, then $[h],[q]\in \cR_{q^*}$.


\item\label{lem:N:d}
If $[h]\in\cL_\alpha$ and $q\cdot h\neq 0$, then $[q\cdot h]\in \cL_\alpha$.\\
If $[h]\in\cR_\alpha$ and $h\cdot q\neq 0$, then $[h\cdot q]\in \cR_\alpha$.
\end{claims}
\end{lemma}

\begin{proof}
\ref{lem:N:a}
The map $\HC\to\HC$ that sends $h$ to $h\cdot\alpha$ is linear
with respect to the underlying vector space~$\langle 1,\qi,\qj,\qk \rangle_\C$
and has kernel
\[
L_\alpha:=\set{h\in\HC}{h\cdot\alpha=0}.
\]
Let
$\beta:=\alpha^*$ and
$V_\beta:=\langle \beta,~\qi\cdot\beta,~\qj\cdot\beta,~\qk\cdot\beta\rangle_\C$.
By assumption,
$
\alpha\cdot\alpha^*=
\beta\cdot\beta^*=
\beta\cdot\alpha=0$ and thus $V_\beta\subseteq L_\alpha$.

Now suppose by contradiction that $\dim V_\beta<2$.
In this case,
\[
\beta=c_1\cdot\qi\cdot\beta=c_2\cdot\qj\cdot\beta=c_3\cdot\qk\cdot\beta
\]
for some $c_1,c_2,c_3\in\C$. Since
\[
\qi=\qj\cdot\qk=-\qk\cdot\qj,\quad
\qj=\qk\cdot\qi=-\qi\cdot\qk,\quad
\qk=\qi\cdot\qj=-\qj\cdot\qi,\quad
\qi^2=\qj^2=\qk^2=-1,
\]
we find that
\[
\begin{array}{rlrlrlrlr}
                \beta &=&          \beta_1 &+&         \beta_2\cdot\qi &+&         \beta_3\cdot\qj &+&         \beta_4\cdot\qk, \\
c_1\cdot\qi\cdot\beta &=& -c_1\cdot\beta_2 &+& c_1\cdot\beta_1\cdot\qi &-& c_1\cdot\beta_4\cdot\qj &+& c_1\cdot\beta_3\cdot\qk, \\
c_2\cdot\qj\cdot\beta &=& -c_2\cdot\beta_3 &+& c_2\cdot\beta_4\cdot\qi &+& c_2\cdot\beta_1\cdot\qj &-& c_2\cdot\beta_2\cdot\qk, \\
c_3\cdot\qk\cdot\beta &=& -c_3\cdot\beta_4 &-& c_3\cdot\beta_3\cdot\qi &+& c_3\cdot\beta_2\cdot\qj &+& c_3\cdot\beta_1\cdot\qk.
\end{array}
\]
By comparing the coefficients, we see that
$\beta_1=-c_i\cdot\beta_{i+1}$,
$\beta_{i+1}=c_i\cdot\beta_1$
and thus
$\beta_1\neq 0$ and
$c_i^2=-1$
for all $1\leq i\leq 3$.
This implies that $\beta_2=\beta_3=\beta_4=\pm\ii\cdot\beta_1$
and thus
\[
\beta\cdot\beta^*=\beta_1^2+\beta_2^2+\beta_3^2+\beta_4^2=-2\cdot\beta_1^2=0.
\]
We arrived at a contradiction as $\beta_1\neq 0$.

We established that $\dim L_\alpha\geq \dim V_\beta\geq 2$ and thus $\dim \cL_\alpha\geq 1$
as $\cL_\alpha=\P(L_\alpha)$ by definition.
Suppose that $h\in L_\alpha$ so that $h\cdot\alpha=0$ and thus $h^*\cdot h\cdot\alpha=h\cdot h^*\cdot\alpha=0$.
Since $h\cdot h^*\in\C$ and $\alpha\neq 0$, we find that $h\cdot h^*=h^*\cdot h=0$.
By definition, $\cN=\P(\set{h\in\HC}{h\cdot h^*=0})$
and thus $\cL_\alpha\subset\cN$.
Since $\cN$ is a smooth hyperquadric of~$\P^3$, it does not contain complex planes.
It follows that $\cL_\alpha$ is a complex line through the complex point~$[\alpha^*]$.
Similarly, we show that $\cR_\alpha\subset\cN$ is a complex line through the complex point~$[\alpha^*]$.
Since $\cL_\alpha\neq\cR_\alpha$, we concluded the proof for Assertion~\ref{lem:N:a}.

For Assertion~\ref{lem:N:b},
notice that if $[h],[q]\in \cL_\alpha$, then
$[h]\in \cL_{q^*}$ by Assertion~\ref{lem:N:a} and thus $h\cdot q^*=0$ by definition.
The remaining Assertions \ref{lem:N:b}, \ref{lem:N:c} and \ref{lem:N:d}
are now a straightforward consequence of the definitions.
\end{proof}

\end{document}
